% analysis/pipeline_story-DE.tex
\section{Pipeline-Erzählung: Von der Messung zum Alignment}

\subsection{Ein einzelner Punkt}
Betrachten wir einen einzelnen gemessenen Punkt \(x\) in der realen Welt. Er
kann aus einer Vermessungskampagne, einem Gleismessfahrzeug oder einem
Entwurfsdatensatz stammen. Auf dieser Stufe existiert der Punkt nur in einem
globalen Koordinatensystem.

\subsection{Schritt 1: CRS-Transformation}
Der Punkt wird zunächst in einem Koordinatenreferenzsystem interpretiert:
\[
x_{\mathrm{crs}}\xrightarrow{\mathcal C}x_{\mathrm{world}}.
\]
Dieser Schritt stellt sicher, dass alle Daten in einem konsistenten räumlichen
Rahmen ausgedrückt sind.

\subsection{Schritt 2: Projektion auf das Gleis}
Der Punkt wird auf das Alignment projiziert:
\[
x_{\mathrm{world}}\xrightarrow{\mathcal W}(s,d).
\]
Dies ist eine nichtlineare Operation: Die longitudinale Position \(s\) wird
bestimmt und die laterale Abweichung \(d\) berechnet. In diesem Augenblick wird
der Punkt Teil des intrinsischen Alignment-Koordinatensystems.

\subsection{Schritt 3: Interpretation durch das Modell}
Das Alignment-Modell stellt eine Krümmungsbeschreibung bereit:
\[
(s,d)\xrightarrow{\mathcal M}\kappa(s).
\]
Dabei wirken parametrische Struktur als Entwurfsabsicht und lokale Korrekturen
als datengetriebene Anpassungen zusammen. Der Punkt wird nicht länger als
isoliertes Datum, sondern als Bestandteil eines strukturierten geometrischen
Systems behandelt.

\subsection{Schritt 4: Physische Realisierung}
Das theoretische Modell wird durch physische Prozesse transformiert:
\[
\kappa(s)\xrightarrow{\mathcal R}\gamma_{\mathrm{real}}(s).
\]
Dieser Schritt berücksichtigt Baubedingungen, mechanisches Verhalten und
Instandhaltungsoperationen. Die entstehende Geometrie ist mit dem theoretischen
Modell nicht identisch.

\subsection{Schritt 5: Zurück zu World-Koordinaten}
Das realisierte Alignment wird in den World-Raum zurückabgebildet:
\[
\gamma_{\mathrm{real}}(s)\xrightarrow{\mathcal W^{-1}}x_{\mathrm{real}}.
\]
Damit entsteht die in der Realität sicht- und messbare Geometrie.

\subsection{Schritt 6: Vergleich und Rückkopplung}
Ursprüngliche Messung und realisierte Position werden verglichen:
\[
x_{\mathrm{real}}\approx x_{\mathrm{world}}.
\]
Die Abweichung treibt die Optimierung: Parameter werden angepasst,
Korrekturen aktualisiert und die Pipeline erneut bewertet.

\subsection{Interpretation}
Der Punkt hat mehrere Darstellungen durchlaufen: globale Koordinaten,
intrinsische Alignment-Koordinaten, parametrisches Modell und realisierte
Geometrie. Jede Transformation führt Struktur, Bedingungen und Interpretation
ein.

\subsection{Zentrale Erkenntnis}
\begin{proposition}
Ein gemessener Punkt ist nicht unmittelbar Teil des Alignments. Er wird erst
durch World--Track-Transformation und Alignment-Modell Teil des Alignments.
\end{proposition}

\subsection{Verbindung zu AIM}
\begin{summary}
Der AIM-Rahmen transformiert räumliche Rohdaten durch eine Operatorfolge in
strukturierte Alignment-Information. Dieser Prozess integriert
Koordinatensysteme, geometrische Modellierung, physische Realisierung und
Optimierung.

Alignment-Modellierung ist damit keine statische Darstellung, sondern ein
dynamischer Prozess, der Messung, Theorie und Bau verbindet.
\end{summary}
